\documentclass{article}
\usepackage{mathtools} 
\usepackage{fontspec}
\usepackage[UTF8]{ctex}
\usepackage{amsthm}
\usepackage{mdframed}
\usepackage{xcolor}
\usepackage{amssymb}
\usepackage{amsmath}


% 定义新的带灰色背景的说明环境 zremark
\newmdtheoremenv[
  backgroundcolor=gray!10,
  % 边框与背景一致，边框线会消失
  linecolor=gray!10
]{zremark}{说明}

% 通用矩阵命令: \flexmatrix{矩阵名}{元素符号}{行数}{列数}
\newcommand{\flexmatrix}[4]{
  \[
  #1 = \begin{pmatrix}
    #2_{11}     & #2_{12}     & \cdots & #2_{1#4}   \\
    #2_{21}     & #2_{22}     & \cdots & #2_{2#4}   \\
    \vdots      & \vdots      & \ddots & \vdots     \\
    #2_{#31}    & #2_{#32}    & \cdots & #2_{#3#4}
  \end{pmatrix}
  \]
}

% 简化版命令（默认矩阵名为A，元素符号为a）: \quickmatrix{行数}{列数}
\newcommand{\quickmatrix}[2]{\flexmatrix{A}{a}{#1}{#2}}

\begin{document}
\title{注释 15.1}
\author{张志聪}
\maketitle

\begin{zremark}
  对按段光滑的理解
\end{zremark}
在分点$x_i$上满足以下条件
\begin{itemize}
  \item 对$f$:
        \begin{align*}
          \lim\limits_{x \to x_i^{-}} f(x), \lim\limits_{x \to x_i^+} f(x) \,\,\, \text{存在}
        \end{align*}

  \item 对$f'$
        \begin{align*}
          \lim\limits_{x \to x_i^{-}} f'(x), \lim\limits_{x \to x_i^+} f'(x) \,\,\, \text{存在}
        \end{align*}
        注：蕴含了左右导数存在
\end{itemize}

\begin{zremark}
  定理15.3 重要性质的证明。
\end{zremark}
\textbf{证明：}

\begin{itemize}
  \item $1^{\circ}$ $f$在$[a, b]$上可积。

        以只有一个分点$x_1$为例.

        我们有
        \begin{align*}
          \lim\limits_{x \to x_1^{-}} f(x)
        \end{align*}
        存在，不妨设为$c$，
        由局部有界性可知，存在$M_1 > 0, U_{-}^{o}(x_1)$，使得
        $x \in U_{-}^{o}(x_1)$有
        \begin{align*}
          |f(x)| < M_1
        \end{align*}

        同理可得，存在$M_2 > 0, U_{+}^{o}(x_1)$，使得
        $x \in U_{+}^{o}(x_1)$有
        \begin{align*}
          |f(x)| < M_2
        \end{align*}
        故存在$x \in U(x_1)$使得
        \begin{align*}
          |f(x)| \leq max\{M_1, M_2, f(x_1)\}
        \end{align*}

        在$U(X_1)$两端各取一个点$x_0, x_2 \in [a, b]$，即$x_0 < x_1 < x_2$，
        由连续函数的有界性可知，在$[a, x_0], [x_2, b]$上，$f$有界。

        综上，$f$在$[a, b]$上有界，于是由定理9.5可知，$f$在$[a, b]$上可积。

\end{itemize}

\end{document}